\section{Poincar\'e 群}

\begin{tikzpicture}[
    scale=0.85,
    transform shape,
    node distance=0.4cm and 0.6cm,
    every node/.style={align=center, font=\sffamily},
    % 基础样式
    base/.style={rectangle, rounded corners=2pt, draw=black!80, thick, inner sep=5pt, fill=white, minimum width=3cm, minimum height=1cm},
    % 颜色区分
    m_color/.style={fill=blue!5, draw=blue!50},
    z_color/.style={fill=red!5, draw=red!50},
    title_style/.style={font=\bfseries\color{gray!80!black}, inner sep=10pt},
    arrow/.style={-stealth, thick, draw=black!50}
]

% --- 核心顶部 ---
\node[base, fill=gray!10, minimum width=2cm] (top) 
    { Poincar\'e 群的不可约表示};

% --- 第一层：Casimir 分支 (错开排布) ---
\node[base, m_color, below left=1.5cm and 0.1cm of top] (m_mass) 
    {\textbf{有质量表示} ($m > 0$)\\ \footnotesize 小群: $SO(3)$};

\node[base, z_color, below right=3.5cm and 0.1cm of top] (z_mass) 
    {\textbf{零质量表示} ($m = 0$)\\ \footnotesize 小群: $ISO(2)$};

% --- 第二层：具体表示 (纵向错位) ---
% 有质量分支
\node[base, m_color, below=1cm of m_mass] (m_0) {$(0,0)$ \\ \footnotesize 标量 (1D)};
\node[base, m_color, below=1cm of m_0] (m_12) {$(\frac{1}{2},0)\oplus(0,\frac{1}{2})$ \\ \footnotesize Dirac 旋量 (4D)};
\node[base, m_color, below=1cm of m_12] (m_1) {$(\frac{1}{2},\frac{1}{2})$ \\ \footnotesize 矢量 (4D)};

% 零质量分支
\node[base, z_color, below=1cm of z_mass] (z_12) {$(\frac{1}{2},0) \text{ or } (0,\frac{1}{2})$ \\ \footnotesize Weyl 旋量 (2D)};
\node[base, z_color, below=1cm of z_12] (z_1) {螺旋度 $h=\pm 1$ \\ \footnotesize 规范矢量 (2 极化)};
\node[base, z_color, below=1cm of z_1] (z_2) {螺旋度 $h=\pm 2$ \\ \footnotesize 规范张量 (2 极化)};

% --- 第三层：场方程 (一一对应) ---
\node[base, m_color, left=1.2cm of m_0] (eq_kg) {Klein-Gordon};
\node[base, m_color, left=1.2cm of m_12] (eq_dirac) {Dirac};
\node[base, m_color, left=1.2cm of m_1] (eq_proca) {Proca};

\node[base, z_color, right=1.2cm of z_12] (eq_weyl) {Weyl};
\node[base, z_color, right=1.2cm of z_1] (eq_maxwell) {Maxwell};
\node[base, z_color, right=1.2cm of z_2] (eq_einstein) {Linearized Einstein};

% --- 连线 ---
\draw[arrow] (top.south) -- ++(0,-0.5) -| (m_mass.north);
\draw[arrow] (top.south) -- ++(0,-0.5) -| (z_mass.north);

% 有质量连线
\draw[arrow] (m_mass) -- (m_0);
\draw[arrow] (m_0) -- (m_12);
\draw[arrow] (m_12) -- (m_1);
\draw[arrow] (m_0.west) -- (eq_kg.east);
\draw[arrow] (m_12.west) -- (eq_dirac.east);
\draw[arrow] (m_1.west) -- (eq_proca.east);

% 零质量连线
\draw[arrow] (z_mass) -- (z_12);
\draw[arrow] (z_12) -- (z_1);
\draw[arrow] (z_1) -- (z_2);
\draw[arrow] (z_12.east) -- (eq_weyl.west);
\draw[arrow] (z_1.east) -- (eq_maxwell.west);
\draw[arrow] (z_2.east) -- (eq_einstein.west);

% --- 补充说明标注 ---
%\node[below=0.8cm of eq_proca, font=\itshape\small, color=blue!70] {* 具有 3 个物理极化};
%\node[below=0.8cm of eq_einstein, font=\itshape\small, color=red!70] {* 具有规范不变性};

\end{tikzpicture}
